\chapter{Related Work}

This research lies at the intersection of Gaussian-based 3D scene reconstruction and reflection--transmission modeling in dynamic environments. We review related research directions: Gaussian Splatting representation, reflection reconstruction, and transparent reconstruction.

\section{Gaussian Splatting}

3D Gaussian Splatting (3DGS)~\cite{kerbl2023} introduces an explicit representation of 3D scenes using a collection of ellipsoid-shaped Gaussian primitives. Each Gaussian is defined by its position, shape, orientation, opacity, and view-dependent color, allowing it to represent a portion of the scene's geometry and appearance. Unlike mesh-based approaches that represent object surfaces using connected polygons, 3DGS uses Gaussian primitives distributed in 3D space to approximate scene geometry and appearance without explicitly constructing continuous surfaces. To generate an image from a given camera viewpoint, 3DGS projects the Gaussian primitives from 3D space onto the 2D image plane. Each projected Gaussian contributes color to a region of the image according to its spatial distribution and opacity. When multiple Gaussians overlap in the image, their contributions are accumulated in depth order, allowing nearby Gaussians to partially or fully occlude those located farther from the camera.

Despite its effectiveness in novel-view synthesis, 3DGS primarily optimizes the observed image appearance without explicitly constraining its volumetric Gaussians to align with actual object surfaces. As a result, the reconstructed geometry may be inaccurate even when the rendered images appear realistic. This limitation is particularly significant for reflection reconstruction where the direction of reflected light depends on the surface normal. Inaccurate surface geometry and orientation can therefore lead to incorrect reflection directions. Furthermore, projecting volumetric Gaussians can produce inconsistent depth estimates for the same surface across different viewpoints, making it complex to determine its precise spatial location. The issue affects transparent scene reconstruction, where reliable depth information is needed to distinguish the glass interface from objects visible behind it.

To address these limitations, 2D Gaussian Splatting (2DGS)~\cite{huang2024} replaces volumetric 3D Gaussian primitives with oriented 2D Gaussian disks embedded in 3D space. Each disk represents a local surface element with a defined position, size, and orientation, allowing the Gaussians to align more closely with actual object surfaces. Furthermore, the orientation of each disk provides a well-defined surface normal, supporting more accurate reconstruction of object surfaces. However, using surface-aligned Gaussian disks alone does not guarantee accurate geometry, since optimizing image appearance can still result in scattered or incorrectly oriented Gaussians. By incorporating depth distortion and normal consistency regularization, 2DGS further improves surface accuracy and geometric consistency across different viewpoints.

The improvements in 2DGS provide a useful geometric foundation for reflection and transparent surface reconstruction. In particular, accurate surface normals are essential for determining reflection directions, while reliable depth and surface geometry help identify the location and orientation of transparent interfaces. Nevertheless, 2DGS does not explicitly model the interaction between reflected and transmitted light.

\section{Reflection Reconstruction}

Reflection reconstruction remains challenging due to the strong view-dependent appearance of reflective surfaces and their complex interactions with surrounding environments. Conventional Gaussian Splatting uses spherical harmonics to represent view-dependent radiance, which often provides insufficient expressiveness for high-frequency specular reflections. To address this limitation, recent studies have explored explicit geometric representations and reflection-aware rendering techniques.

MirrorGaussian~\cite{liu2024} reconstructs mirror reflections by exploiting the geometric symmetry between the real scene and its mirror image, using 3D Gaussians to represent the scene and an estimated mirror plane to determine the reflection geometry. By reflecting the original Gaussians across this plane, the method constructs a mirrored scene in which the Gaussian primitives are positioned symmetrically on the opposite side of the mirror, enabling high-quality reflection reconstruction without explicitly tracing reflected rays. In contrast, EnvGS~\cite{xie2025} reconstructs reflections through an explicit representation of the surrounding environment rather than relying on geometric symmetry. Specifically, it employs base Gaussians to represent scene geometry and appearance, alongside a separate set of environment Gaussians to model the surrounding content that contributes to reflections. By determining reflection directions from surface geometry and tracing rays into the environment Gaussians, EnvGS retrieves the reflected colors corresponding to the surrounding objects. This approach enables the reconstruction of detailed reflections from nearby objects, which are difficult to reproduce using conventional environment maps that do not explicitly represent their spatial positions.

Recent studies have further extended reflection reconstruction to dynamic environments. SpectroMotion~\cite{fan2025} combines deformation fields with physically based rendering to reconstruct moving reflective objects and their time-varying appearances. Furthermore, it introduces a deformable environment map to represent changes in reflected appearance over time. However, its environment map representation remains limited in capturing detailed reflections from nearby objects. To overcome this limitation, TraceFlow~\cite{tao2025} represents the surrounding environment using a separate set of dynamic Gaussians with explicit spatial information to capture near-field reflections. While these methods excel at representing high-frequency specular effects, they still primarily focus on reflective surfaces and do not explicitly account for the coexistence of reflection and transmission through transparent surfaces.

\section{Transparent Scene Reconstruction}

Unlike purely reflective surfaces, transparent surfaces allow both reflected and transmitted light to contribute to the observed appearance. Recent advances in Gaussian-based transparent scene reconstruction have focused on recovering accurate surface geometry and explicitly modeling the reflection and transmission of light through transparent materials. In conventional methods, reducing the opacity of transparent surfaces allows background objects to remain visible but also reduces the contribution of these surfaces to depth estimation, resulting in inaccurate reconstructed geometry. To address this transparency-depth dilemma, TSGS~\cite{li2025} separates geometry reconstruction from appearance refinement, learning surface geometry with reduced specular interference before improving visual appearance while preserving the reconstructed geometry. Beyond geometric accuracy, realistic transparent reconstruction also requires modeling how surrounding objects contribute to the observed appearance through reflection and light transmission. TransparentGS~\cite{huang2025} combines transparent Gaussian primitives with Gaussian light field probes that store color and depth information about the surrounding environment. By querying these probes, the method captures reflections and the appearance of objects visible through transparent materials. Nevertheless, these approaches primarily focus on the geometry and appearance of transparent objects, rather than reconstructing reflected and transmitted content as independent 3D scenes across planar glass surfaces.

For glass surfaces where reflection and transmission overlap, reconstructing the observed appearance requires accounting for both the scene visible through the glass and the surrounding scene reflected on its surface. TR-Gaussians~\cite{liu2025} represents the real scene using primary Gaussians, which are used to render the transmission image. To generate reflections, the method mirrors these Gaussians across a learnable glass plane, creating a virtual scene whose rendering produces the reflection image. The reflected and transmitted images are then combined according to the viewing direction. In contrast, RT-Splatting~\cite{shi2026} introduces a unified Gaussian representation that separates geometric occupancy from optical opacity. Geometric occupancy determines the contribution of Gaussians to surface reconstruction, while optical opacity controls light attenuation to preserve the transmitted background. This separation enables sharp reflections and clear transmission within a shared representation.

Extending reflection--transmission reconstruction to scene-scale transparency, GLINT \cite{na2026} explicitly separates the scene into three Gaussian components: interface, transmission, and reflection. The interface Gaussians reconstruct the first visible surface and provide geometric and material information to guide the rendering of the other two components. GLINT combines interface reconstruction with ray tracing to retrieve reflected and transmitted radiance from their respective Gaussian representations. This enables consistent reconstruction of transparent surfaces and the scene content visible through and reflected by them.

Despite these advances, existing transparent scene reconstruction methods primarily focus on static environments. Although reflection and transmission can be explicitly represented, these methods do not directly model scenes where the reflected and transmitted content changes independently over time. Extending this decomposition to dynamic environments therefore remains an open research direction.
